\documentclass[11pt]{article}
\setlength{\textwidth}{6.5in}
\setlength{\oddsidemargin}{0in}
\setlength{\evensidemargin}{0in}
\setlength{\textheight}{9in}
\setlength{\topmargin}{-0.5in}
\usepackage{amsmath,amssymb,amsthm}
\usepackage[hidelinks]{hyperref}
\newtheorem{theorem}{Theorem}
\newtheorem{proposition}[theorem]{Proposition}
\newtheorem{example}[theorem]{Counterexample}
\newcommand{\E}{\mathbb E}
\newcommand{\R}{\mathbb R}
\newcommand{\eps}{\varepsilon}
\newcommand{\Cov}{\operatorname{Cov}}
\newcommand{\Var}{\operatorname{Var}}
\newcommand{\dist}{\operatorname{dist}}
\title{Proof addendum and mathematical corrections\\to the conditional-flow manuscript}
\author{Independent audit of the local manuscript}
\date{16 October 2026}
\begin{document}
\maketitle
This addendum is separate from \texttt{paper/main.tex}. Its source revision and
review ledger are recorded in \texttt{06\_Math\_Formalization.md}. The first six
results each received four fresh adversarial confirmations under their stated
assumptions. Subsequent corrections have explicit proofs but have not all received
that review depth. None of the results identifies a trained neural network with
an exact population field.

\section{Hard conditioning under guidance}
Guidance need not preserve the full trajectory. It nevertheless preserves the
conditioned endpoint without assuming convergence or excluding nearest-atom ties.
\begin{theorem}\label{thm:hard}
Let $a_1,\ldots,a_N\in\R^d$, $N\ge1$, fix $i$, and put $\eps=1-t$ and
\[
q_j(x,t)=\frac{\exp(-\lVert x-ta_j\rVert^2/(2\eps^2))}
 {\sum_\ell\exp(-\lVert x-ta_\ell\rVert^2/(2\eps^2))},
\qquad m(x,t)=\sum_jq_j(x,t)a_j.
\]
For every $w\ge0$ and $x_0\in\R^d$, the solution of
\[
\dot x_t=\frac{(1+w)a_i-wm(x_t,t)-x_t}{1-t},\qquad x_{t=0}=x_0,
\]
exists on $[0,1)$ and converges to $a_i$. More precisely, for $0<t_0<t<1$,
\begin{equation}\label{eq:hard-bound}
\lVert x_t-ta_i\rVert^2\le\eps^2\left[
\frac{\lVert x_{t_0}-t_0a_i\rVert^2}{(1-t_0)^2}
+2w\log N\log\frac{t(1-t_0)}{t_0(1-t)}\right].
\end{equation}
\end{theorem}
\begin{proof}
On every compact time interval below $1$ the field is smooth, locally Lipschitz,
and has linear growth, since $m$ lies in the convex hull of the finite atoms.
These properties give existence and uniqueness on $[0,1)$.

Write $D_j=\lVert x-ta_j\rVert^2$ and $H(q)=-\sum_jq_j\log q_j$.
The Gibbs identity and the contribution of index $i$ to the partition function give
\[
\sum_jq_jD_j=2\eps^2\bigl(H(q)-\log\sum_je^{-D_j/(2\eps^2)}\bigr)
\le D_i+2\eps^2\log N.
\]
With $z=x-ta_i$, expanding $D_j-D_i$ therefore gives
\[
2t z\cdot(a_i-m)+t^2\sum_jq_j\lVert a_i-a_j\rVert^2
\le2\eps^2\log N.
\]
Along the flow, $\dot z=(-z+w(a_i-m))/\eps$. Dropping the nonnegative
sum for $t>0$ and using $w\ge0$, $Y=\lVert z\rVert^2$ satisfies
\[
Y'\le-\frac{2Y}{\eps}+\frac{2w\eps\log N}{t},\qquad
\left(\frac{Y}{\eps^2}\right)'\le\frac{2w\log N}{t(1-t)}.
\]
Integration yields \eqref{eq:hard-bound}. Its right side tends to zero, so
$x_t-ta_i\to0$ and $x_t\to a_i$.
\end{proof}
The result applies to the exact hard-conditioned population field. It asserts
neither convergence of a learned field nor atom selection for positive label
bandwidth. For $w=-1$, symmetric atoms $\{-1,1\}$ admit the unconditional
trajectory $x_t=0$, so the sign restriction cannot simply be dropped.

\section{The sharp Lipschitz floor}
\begin{theorem}\label{thm:sharp}
Let $I$ be uniform on finitely many indices, let $T$ be uniform on $(0,1)$,
and let $X_0\sim\mathcal N(0,I_d)$, mutually independently. The distinct
conditioning labels $y_i$ identify atoms $a_i$. Set
$X_T=(1-T)X_0+Ta_I$ and $U=a_I-X_0$.
For the class $\mathcal F_L$ of measurable fields that are spatially
$L$-Lipschitz at every time and supplied label, $L\ge1$,
\[
\inf_{v\in\mathcal F_L}\E\lVert v(X_T,T,y_I)-U\rVert^2=\frac{d}{3L}.
\]
An optimizer, for $t<1$, is
\[
v_L(x,t,y_i)=a_i-\min\{L,(1-t)^{-1}\}(x-ta_i).
\]
\end{theorem}
\begin{proof}
Fix $i,t$ and write $\eps=1-t$. In this paragraph expectations are conditional
on $I=i,T=t$, so $X_t=ta_i+\eps X_0$; abbreviate $v(X_t,t,y_i)$ by $v(X_t)$.
For an independent copy $X_t'$ of $X_t$,
\[
\E\lVert v(X_t)-\E v(X_t)\rVert^2
=\tfrac12\E\lVert v(X_t)-v(X_t')\rVert^2\le dL^2\eps^2.
\]
The centered target is $U-\E U=-X_0$, with $L^2$ norm $\sqrt d$. Centering the residual and
using the reverse triangle inequality gives conditional risk at least
$d[1-L\eps]_+^2$. Infinite-risk fields satisfy this bound trivially.
For the displayed optimizer the residual is
$[1-\min\{L,\eps^{-1}\}\eps]X_0$, which attains that bound at every $i,t$.
Its sampled value has uniformly bounded second moment. Finally,
\[
d\int_0^1[1-L(1-t)]_+^2\,dt=\frac{d}{3L}.
\]
\end{proof}
This optimizes over all measurable spatially Lipschitz fields, rather than a
particular architecture. A small lower bound alone does not upper-bound a
network's approximation error or rule out a representation gap.

\section{A spectral obstruction to a probability kernel}
\begin{theorem}\label{thm:filter}
For $a,\zeta>0$, the continuous function
\[
G(0)=1,\qquad G(\omega)=1-e^{-a|\omega|^{-\zeta}}\quad(\omega\ne0)
\]
is not the characteristic function of a probability measure on $\R$.
\end{theorem}
\begin{proof}
If $G(\omega)=\int e^{i\omega x}\,d\mu(x)$ for a probability measure $\mu$,
fix a sequence $\omega_n\ne0$ tending to zero. Fatou's lemma gives
\[
\frac12\int x^2\,d\mu(x)
\le\liminf_{n\to\infty}\int\frac{1-\cos(\omega_n x)}{\omega_n^2}\,d\mu(x)
=\lim_{n\to\infty}\frac{e^{-a|\omega_n|^{-\zeta}}}{\omega_n^2}=0.
\]
Thus $\mu=\delta_0$, whose characteristic function is identically $1$,
contradicting $G(\omega)<1$ for $\omega\ne0$.
\end{proof}
No moment assumption on $\mu$ is needed. Hence this continuum multiplier cannot
represent convolution by a probability measure for all targets. Signed kernels,
distributional kernels, and approximate fits on a finite set of conditions are
not excluded. In particular, a good moment fit to a Gaussian-label mixture is
not an identity of distributions.

\section{Convergence after averaging over labels}
The conditional convergence statement for an averaged-label field can be
strengthened. The argument below proves convergence, while retaining a genuine
distinction between atom limits and ties.
\begin{theorem}\label{thm:averaged}
Let $a_1,\ldots,a_N\in\R^d$, let $\phi$ be a fixed probability measure,
and let $c_j(y)$ be measurable with $0<c_j(y)<\infty$ almost everywhere.
Set $\eps=1-t$ and
\[
q_j(x,t,y)=\frac{c_j(y)e^{-\lVert x-ta_j\rVert^2/(2\eps^2)}}
 {\sum_\ell c_\ell(y)e^{-\lVert x-ta_\ell\rVert^2/(2\eps^2)}},\qquad
m(x,t)=\int\sum_jq_j(x,t,y)a_j\,d\phi(y).
\]
For every $x(0)=x_0\in\R^d$, the equation $x'=(m(x,t)-x)/(1-t)$ has a unique
solution on $[0,1)$. It converges to a point $z\in\operatorname{conv}N(z)$, where
$N(z)=\{a_j:\lVert z-a_j\rVert=\min_\ell\lVert z-a_\ell\rVert\}$.
No integrability assumption on logarithms of the coefficients
is required. A non-atom limit is possible.
\end{theorem}
\begin{proof}
The original field is locally Lipschitz on compact intervals $t<1$, including
$t=0$: its spatial softmax derivatives are bounded independently of the
coefficients there. Since $m$ is in the convex hull of the atoms,
variation of constants gives existence on all of $[0,1)$ and a bounded trajectory.
We now prove convergence of this trajectory.

\emph{A decreasing energy.}
Normalize $p_j(y)=c_j(y)/\sum_\ell c_\ell(y)$. Start at any $t_0>0$ and set
\[
\tau=-\log(1-t),\quad \alpha=\frac{\eps^2}{t},\quad
b_j=x\cdot a_j-\frac t2\lVert a_j\rVert^2,\quad R=\max_j\lVert a_j\rVert.
\]
From now on dots denote differentiation with respect to $\tau$, with
$t=1-e^{-\tau}$. Then $q_j\propto p_je^{b_j/\alpha}$ and $\dot x=m-x$,
which also gives a bounded $\tau$-speed. Define
\[
F=\int\alpha\log\sum_jp_j(y)e^{b_j/\alpha}\,d\phi(y),\qquad
W=\tfrac12\lVert x\rVert^2-F+\tfrac12\eps R^2.
\]
The integrand lies between $\min b_j$ and $\max b_j$. Treating $t$ and $\alpha$
as independent arguments, differentiation gives
\[
\nabla_xF=m,\qquad \partial_\alpha F=-\int D_{\rm KL}(q\Vert p)\,d\phi,
\qquad \partial_tF=-\tfrac12\int\sum_jq_j\lVert a_j\rVert^2\,d\phi.
\]
These derivatives pass under the integral on compact positive-time intervals:
the spatial and time derivatives are bounded, and
\[
0\le D_{\rm KL}(q\Vert p)
=\frac{\sum_jq_jb_j}{\alpha}-\log\sum_jp_je^{b_j/\alpha}
\le\frac{\max_jb_j-\min_jb_j}{\alpha}
\]
is uniform in $y$. Since $\dot\alpha<0$ and $\dot t=\eps$,
\[
\dot W=-\lVert\dot x\rVert^2+\dot\alpha\int D_{\rm KL}(q\Vert p)\,d\phi
+\frac\eps2\left(\int\sum_jq_j\lVert a_j\rVert^2\,d\phi-R^2\right)
\le-\lVert\dot x\rVert^2.
\]
Moreover $F\le\max_jb_j$, so
\[
W\ge\min_j\left\{\tfrac12\lVert x-a_j\rVert^2
+\tfrac\eps2(R^2-\lVert a_j\rVert^2)\right\}\ge0.
\]
Thus $\int_{\tau_0}^\infty\lVert\dot x\rVert^2\,d\tau<\infty$.

\emph{A finite set of possible limits.}
If $x_n\to z$ and $\tau_n\to\infty$, every weight on an atom outside $N(z)$
tends to zero for each $y$, by the positive distance gap and the fixed positive
coefficient ratios. Bounded convergence gives the same conclusion after
averaging over $y$. Every limit of $m(x_n,t_n)$ therefore lies in
$\operatorname{conv}N(z)$.

Let $C=\{z:z\in\operatorname{conv}N(z)\}$. This set is finite. For each
nonempty subset $S$ of atoms, there is at most one point in
$\operatorname{aff}S$ equidistant from all atoms of $S$: the difference of two
such points lies both in $\operatorname{span}(S-S)$ and its orthogonal
complement. A point of $C$ is one of these finitely many candidates.
On any compact set disjoint from $C$, the speed is eventually bounded below
by a positive constant. Otherwise a sequence of points there with speed tending
to zero would, by compactness and the preceding concentration argument, have
a limit in $C$.

An omega-limit point outside $C$ admits a closed ball disjoint from $C$.
Infinitely many visits to its inner half force residence for a fixed positive
time in the ball, because the speed is bounded above. If the trajectory remains
in the ball forever, the eventual positive lower speed bound already contradicts
finite energy. Otherwise choose disjoint visits; each spends a fixed positive
amount of energy, again a contradiction. Hence the omega-limit set is contained
in $C$. It is nonempty, compact, and connected, as the intersection of the nested
closures of the bounded connected trajectory tails. Being contained in a finite
set, it is a singleton.

Finally, equal coefficients and atoms $\{-a,a\}$ give the trajectory $x_t=0$
from $x_0=0$. Its limit is not an atom when $a\ne0$.
\end{proof}
This theorem applies to a positive-bandwidth label average in the manuscript.
Its exceptional tie limits do not preclude an almost-sure atomic endpoint.
The next result addresses the measure of their basins.

\section{Atomic endpoints almost surely}
The statements in this section concern exact fields with fixed conditioning
parameters. Combine repeated atom locations before applying the first theorem;
the second theorem is stated for distinct atoms. Throughout this section
$\tau=-\log(1-t)$, $\alpha=(1-t)^2/t$, and $\beta=\alpha^{-1}$; dots denote
$\tau$ derivatives. Energy arguments start at any $\tau_0>0$.
\begin{theorem}\label{thm:average-as}
Under the assumptions of Theorem~\ref{thm:averaged}, if the initial law is
absolutely continuous with respect to Lebesgue measure, its endpoint law is
supported on the finite atom set.
\end{theorem}
\begin{proof}
We may assume the atoms distinct, summing coefficients at repeated locations.
Use the $\tau$ variables of the preceding proof and write $\beta=1/\alpha$.
For a point $r$ in the closed probability simplex, define
\[
Q_j(y,r)=\frac{p_j(y)r_j}{\sum_\ell p_\ell(y)r_\ell},\quad
M(r)=\int\sum_jQ_j(y,r)a_j\,d\phi,\quad
V(r)=\int\operatorname{tr}\Cov_{Q(y,r)}(a)\,d\phi.
\]
All denominators are positive for each $y$ outside one fixed null set.
The integrands are continuous in $r$ and uniformly bounded. Dominated
convergence makes $M,V$ continuous on the closed simplex. Positivity of the
coefficients and distinctness of the atoms imply that $V(r)=0$ exactly at its
vertices. Compactness therefore gives, for each $\delta>0$, a constant
$\kappa_\delta>0$ such that
\[
\dist(M(r),A)\ge\delta\quad\Longrightarrow\quad V(r)\ge\kappa_\delta,
\qquad A=\{a_1,\ldots,a_N\},
\]
whenever the set on the left is nonempty.
For the actual unconditional softmax vector $r_j\propto e^{b_j/\alpha}$,
the averaged $\tau$-field $f=m-x$ has divergence
\[
\operatorname{div}f=\beta V(r)-d.
\]

Consider a trajectory converging to a non-atom $z$, with $h=\dist(z,A)>0$.
Eventually $\lVert x-z\rVert<h/4$. At good times
$\lVert\dot x\rVert\le h/4$, the identity $m=x+\dot x$ gives
$\dist(m,A)\ge h/2$, hence $V\ge\kappa_{h/2}$.
The complementary times have finite measure by the finite energy proved above.
Their measure in the last unit interval therefore tends to zero. Good times
thus occupy at least half of every sufficiently late unit interval. Since
$V\ge0$ everywhere and $\beta(\tau)=(1-e^{-\tau})e^{2\tau}$,
\[
\int_{\tau_0}^\tau\operatorname{div}f(x(s),s)\,ds
\ge\kappa_{h/2}\int_{\mathrm{good}\cap[\tau-1,\tau]}\beta(s)\,ds
-d(\tau-\tau_0)\longrightarrow+\infty.
\]

The finite-time flow maps $F_\tau=F_{\tau,\tau_0}$ from $\tau_0>0$ are $C^1$
diffeomorphisms: spatial derivatives are locally uniformly bounded, and linear
growth permits both forward and backward solutions on finite intervals.
Liouville's formula gives
\[
J_\tau(\xi):=\det DF_\tau(\xi)
=\exp\!\left(\int_{\tau_0}^\tau
\operatorname{div}f(F_s(\xi),s)\,ds\right)>0.
\]
Thus $J_\tau\to\infty$ at every non-atom-convergent initial point.
For a bounded initial ball $K$, boundedness of $m$ places all $F_\tau(K)$
in one fixed ball $B$. Let $E\subset K$ be the non-atom basin; it is measurable,
since the limit map is a pointwise limit of continuous flow maps. Injectivity
and change of variables give $\int_EJ_\tau\,dx\le\operatorname{Leb}(B)$.
Fatou's lemma along integer times forces $\operatorname{Leb}(E)=0$.
Exhaust the initial space by bounded balls and take the preimage of the null
set under the finite-time diffeomorphism from $t=0$ to $t(\tau_0)$. Its inverse
is locally Lipschitz and preserves Lebesgue-null sets. Absolute continuity
of the source proves the assertion.
\end{proof}

\begin{theorem}\label{thm:cfg-as}
Let the finite atoms $a_j$ be distinct, let $p_c$ be a fixed strictly positive
probability vector, let $p_u$ be uniform, and fix a finite $w\ge0$.
Define $q_{p,j}\propto p_j\exp(-\lVert x-ta_j\rVert^2/(2(1-t)^2))$,
normalized to sum to one, and $m_p=\sum_jq_{p,j}a_j$.
Write $m_c=m_{p_c}$ and $m_u=m_{p_u}$; this formula also defines the field at
$t=0$. Every trajectory from a finite initial point at $t=0$ of
\[
\frac{dx}{d\tau}=(1+w)m_c(x,t)-wm_u(x,t)-x
\]
converges. Its limit is a point $z\in\operatorname{aff}N(z)$. For an absolutely
continuous initial law, the endpoint is a training atom almost surely.
\end{theorem}
\begin{proof}
\emph{Convergence.}
Put $H_p=\alpha\log\sum_jp_je^{b_j/\alpha}$ and
\[
U=\tfrac12\lVert x\rVert^2-(1+w)H_{p_c}+wH_{p_u}.
\]
The mean $(1+w)m_c-wm_u$ is bounded, so trajectories and $\tau$-speeds
are bounded and finite-time flows are well posed. We have $\nabla_xU=-\dot x$.
For fixed positive $p$, $D_{\rm KL}(q_p\Vert p)\le\log(1/\min_jp_j)$.
The derivative identities above therefore give
\[
\dot U=-\lVert\dot x\rVert^2+\mathcal R(\tau),\qquad
|\mathcal R(\tau)|\le C(|\dot\alpha|+\dot t).
\]
The latter bound is integrable. The min--max bounds on $H_p$ show that $U$
is bounded below on a bounded trajectory. Hence its squared speed is integrable.

Near any cluster point $z$, the non-nearest softmax weights tend uniformly to
zero at large times. Thus the signed mean approaches $\operatorname{aff}N(z)$.
If $z$ were outside that affine hull, its orthogonal distance would force a
positive lower speed bound in a fixed neighborhood at late times. Bounded upper
speed and recurrent visits would then contradict finite energy, as in
Theorem~\ref{thm:averaged}. Each cluster point consequently lies in
$\operatorname{aff}N(z)$. There are finitely many such equidistant candidates,
and connectedness of the omega-limit set proves convergence.

\emph{The non-atom basins have zero volume.}
Consider a limit $z\notin A$, and let $r>0$ be the common distance of its nearest
atoms. Write $\mu_c=m_c-z$, $\mu_u=m_u-z$, and
$\mu_w=(1+w)\mu_c-w\mu_u$. Write $\Cov_c,\Cov_u$ for covariance under
$q_{p_c},q_{p_u}$, respectively. The signed covariance trace is
\begin{align*}
T&=(1+w)\operatorname{tr}\Cov_c(a)-w\operatorname{tr}\Cov_u(a)\\
 &=r^2-(1+w)\lVert\mu_c\rVert^2+w\lVert\mu_u\rVert^2+o(1)\\
 &=r^2-\lVert\mu_w\rVert^2+
 \frac{w}{1+w}\lVert\mu_u-\mu_w\rVert^2+o(1)
 \ge r^2-\lVert\mu_w\rVert^2+o(1).
\end{align*}
The error tends to zero along the whole trajectory, since non-nearest mass
vanishes and the atoms are finite. Since $\mu_w=\dot x+x-z$, convergence and
finite energy give $\int_n^{n+1}\lVert\mu_w\rVert^2\,d\tau\to0$.
On each late unit interval the ratio $\beta_{\max}/\beta_{\min}$ is bounded.
Writing the preceding error as at worst $-\eta_n$ throughout that interval,
where $\eta_n\to0$, yields
\[
\int_n^{n+1}\beta T\,d\tau
\ge(r^2-\eta_n)\beta(n)-\beta(n+1)
\int_n^{n+1}\lVert\mu_w\rVert^2\,d\tau
\ge\tfrac14 r^2\beta(n)
\]
eventually. The divergence is $\beta T-d$, so its accumulated integral tends
to infinity along integer times. The same bounded-image-ball,
change-of-variables, and Fatou argument as in Theorem~\ref{thm:average-as}
makes every non-atom basin null. Absolute continuity completes the proof.
\end{proof}
For a Gaussian label kernel at any fixed positive bandwidth and fixed query,
$p_c$ is strictly positive, so this theorem closes the endpoint premise of
the manuscript's positive-bandwidth guidance statement. The endpoint weights
are not identified, and the theorem asserts nothing about errors of learned
fields. Symmetric tie trajectories remain possible for exceptional initial
points.

\section{Corrections to statistical statements}
\subsection{A population smoother is not a finite-sample ratio}
Let $\mu(y)=By$, $Y\sim\mathcal N(0,\Sigma_Y)$, and use a Gaussian label kernel.
The deterministic smoothed population mean is
\[
\mu_h^\infty(y):=\frac{\E[K_h(y-Y)X]}{\E[K_h(y-Y)]}
=B\Sigma_Y(\Sigma_Y+h^2I)^{-1}y.
\]
Indeed the reweighted label law is Gaussian, by completing its quadratic form,
with mean $\Sigma_Y(\Sigma_Y+h^2I)^{-1}y$. Consequently
\[
\mu_h^\infty(y)-\mu(y)=-Bh^2(\Sigma_Y+h^2I)^{-1}y.
\]
The empirical ratio $\widehat\mu_h=\sum K_hX_i/\sum K_h$ is random and is not
equal to this expression. Under the manuscript's local regularity assumptions,
the appropriate pointwise statement is
\[
\widehat\mu_h-\mu=(\mu_h^\infty-\mu)+O_P((Nh^k)^{-1/2}).
\]
Also $B\Sigma_Y^{-1}y\ne0$ for every $y\ne0$ requires an injectivity
assumption on $B$; it does not follow for arbitrary dimensions and rank.

\begin{example}[Local assumptions do not imply finite expected risk]
Choose a nonnegative label $Y$ with a bounded density, smooth and positive on
a neighborhood of $y=1$, and with $\Pr(Y>u)\sim1/\log u$ at infinity.
Let $X=Y+Z$, where $Z$ is independent uniform on $[0,1]$.
Then $\mu(y)=y+1/2$, $\Sigma(y)=1/12$, and the conditional fourth moment has
polynomial growth. These satisfy the local random-design assumptions.
Nevertheless every positive-kernel empirical mean obeys
\[
\widehat\mu_h(1)\ge\min_{i\le N}Y_i,\qquad
\E\min_{i\le N}Y_i=\int_0^\infty\Pr(Y>u)^N\,du=\infty.
\]
Thus finite expected bias and mean squared error need additional tail or
uniform-integrability assumptions; the $O_P$ statement is not contradicted.
Such a density exists by smoothly joining a bounded positive density near $1$
to a constant multiple of $1/(u(\log u)^2)$ beyond a fixed large threshold.
\end{example}

\subsection{Actual label-conditioned weighted moments}
Condition on all labels and suppose $X_i$ are independent with means $\mu_i$
and covariances $\Sigma_i$. For weights $p_i$ determined by those labels, let
$\bar\mu=\sum p_i\mu_i$ and $C_p=\sum p_i(X_i-\sum p_jX_j)(X_i-\sum p_jX_j)^\top$.
Expanding $C_p=\sum p_iX_iX_i^\top-\bar X\bar X^\top$ gives exactly
\begin{align*}
\Cov(\bar X\mid Y)&=\sum_i p_i^2\Sigma_i,\\
\E[C_p\mid Y]&=\sum_i p_i(1-p_i)\Sigma_i
+\sum_i p_i(\mu_i-\bar\mu)(\mu_i-\bar\mu)^\top.
\end{align*}
The simpler formulas $S/n_{\mathrm{eff}}$ and $(1-1/n_{\mathrm{eff}})S$ require
common conditional means and covariance $S$, or fixed weights independent of
i.i.d. atoms. Applying them to actual local atoms is a model approximation.

\subsection{Selecting bandwidth using responses}
\begin{proposition}
Suppose the manuscript's label-mixture coverage theorem applies separately at
every bandwidth in a fixed deterministic grid $\mathcal H$ of size $G\ge1$.
Then, with probability at least $1-\delta$, simultaneously for every
$h\in\mathcal H$ and every nominal coverage $\alpha$,
\[
|C_{\alpha,h}-\alpha|\le 2L\E r_h+8\E n_{\mathrm{eff},h}^{-1/2}
+\sqrt{2a_{\max,h}\log(G/\delta)}.
\]
The same bound holds at any measurable selected grid value
$\widehat h(X_1,Y_1,\ldots,X_N,Y_N)$.
\end{proposition}
\begin{proof}
Conditional on the labels, apply the original theorem at each fixed $h$ with
failure probability $\delta/G$. The union bound gives a single event for the
entire grid, including any selected element. Integrate over the labels.
\end{proof}
Directly substituting a response-fitted bandwidth into the original fixed-weight
theorem omits this selection step. A continuous bandwidth search needs a separate
uniform argument. The CDF continuity condition for a two-sided bound remains.

\subsection{Variance estimation needs a fourth moment}
For $M\ge2$ i.i.d. vectors, mean $\mu$, covariance $\Sigma$, and finite fourth
moment, let $T$ be the trace of the covariance with divisor $M-1$. Then
\begin{equation}\label{eq:vartrace}
\Var T=\frac{\E\lVert X-\mu\rVert^4-(\operatorname{tr}\Sigma)^2}{M}
+\frac{2\operatorname{tr}(\Sigma^2)}{M(M-1)}.
\end{equation}
To verify it, write $T$ as the order-two U-statistic with kernel
$g(X,X')=\lVert X-X'\rVert^2/2$. Pairs sharing one index have covariance
$[\E\lVert X-\mu\rVert^4-(\operatorname{tr}\Sigma)^2]/4$;
$\Var g$ is half that fourth-moment difference plus $\operatorname{tr}(\Sigma^2)$.
Counting overlapping and equal pairs yields \eqref{eq:vartrace}.
For a Gaussian law the first numerator is $2\operatorname{tr}(\Sigma^2)$, so,
when $\operatorname{tr}\Sigma>0$,
\[
\frac{\sqrt{\Var T}}{\operatorname{tr}\Sigma}
=\sqrt{\frac{2\operatorname{tr}(\Sigma^2)}{(M-1)(\operatorname{tr}\Sigma)^2}}
\le\sqrt{\frac2{M-1}}.
\]
This is not distribution-free: in one dimension Bernoulli$(p)$ has kurtosis
$(1-3p+3p^2)/(p(1-p))$, which diverges as $p\downarrow0$.

\subsection{Exchangeable ranks with ties}
Let $S_0,\ldots,S_M$ be exchangeable scores, with $S_0$ the truth's score. Put
$A=\#\{j\ge1:S_j<S_0\}$ and $B=\#\{j\ge1:S_j=S_0\}$.
Without ties, $(A+1/2)/(M+1)$ is uniform on the $M+1$ midpoint ranks. With ties,
the deterministic midrank $(A+B/2+1/2)/(M+1)$ need not be uniform: identical
scores always give $1/2$.

To recover the discrete uniform rank, independently choose
$J\mid S\sim\operatorname{Unif}\{0,\ldots,B\}$ and use
$(A+J+1/2)/(M+1)$. Randomly ordering each tied block makes each exchangeable
index equally likely to occupy every rank, proving the assertion.
Alternatively, with independent $V\sim\operatorname{Unif}(0,1)$,
$[A+V(B+1)]/(M+1)$ is continuous uniform. For joint scores computed from the
sample, the score construction itself must be permutation-equivariant across
truth and draws, for example symmetric leave-one-out scores with a common
ridge rule. Correct marginal distributions alone do not imply exchangeability.

\section{What endpoint smoothing can improve}
The nearest-atom obstruction is a lower bound for every atom-supported
approximation. It is not an upper bound on the benefit of adding noise to a
particular, poorly weighted approximation.
\begin{example}\label{ex:smoothing}
Let $A=\{-1,1\}$, $P_0=(\delta_{-1}+\delta_1)/2$, and $Q_0=\delta_{-1}$.
Then $F(P_0):=\int\dist(x,A)^2\,dP_0(x)=0$. For $Q_\rho=\mathcal N(-1,\rho^2)$,
\[
W_2^2(Q_\rho,P_0)=2+\rho^2-2\rho\sqrt{2/\pi}.
\]
In fact the increasing transport sends the lower Gaussian half to $-1$ and
the upper half to $1$. With $Z$ standard normal its cost is
$\E[\rho^2Z^2\mathbf1_{Z\le0}+(\rho Z-2)^2\mathbf1_{Z>0}]$,
which is the displayed expression. At $\rho=\sqrt{2/\pi}$, the improvement
over $W_2^2(Q_0,P_0)=2$ is $2/\pi>F(P_0)$.

This persists for a continuous target and positive weights on both atoms.
Set $P_\tau=P_0*\mathcal N(0,\tau^2)$ and
$Q_\eta=(1-\eta)\delta_{-1}+\eta\delta_1$. Coupling gives
$F(P_\tau)\le\tau^2$, $W_2(P_\tau,P_0)\le\tau$, and
$W_2(Q_\eta,Q_0)\le2\sqrt\eta$. Gaussian convolution contracts $W_2$ by
using a common noise. The triangle inequality therefore makes both costs
continuous as $\eta,\tau\downarrow0$, and the strict improvement remains
larger than $F(P_\tau)$ for sufficiently small positive $\eta,\tau$.
\end{example}
For a target with mean $\mu_*$ and covariance $\Sigma_*$, a Gaussian-noised
atom mixture can match both moments exactly only if some $h$ satisfies
$\bar a_h=\mu_*$ and $\Sigma_*-C_h=\rho^2I_d$ for a scalar $\rho^2\ge0$.
Two scalar controls do not guarantee these vector and matrix equations.
\end{document}
